% !TEX root = ../main.tex
\begin{abstractind}
  \justifying
  Mobilitas kendaraan roda dua yang meningkat di lingkungan perguruan tinggi mempertinggi risiko keselamatan, terutama akibat ketidakpatuhan penggunaan helm. Pengawasan manual yang berlaku saat ini kurang terukur dan belum memberi wawasan spasiotemporal yang presisi sebagai dasar evaluasi kebijakan. Penelitian ini merancang dan mengevaluasi arsitektur pengawasan cerdas berbasis aliran video waktu nyata (\textit{real-time video streaming}) secara ujung ke ujung (\textit{end-to-end}). Arsitektur tersebut memadukan YOLOv8 untuk ekstraksi fitur spasial, ByteTrack untuk pelacakan multiobjek, filter \textit{Intersection over Area} (IoA) dan konfirmasi temporal \textit{N-Frame} untuk menekan deteksi palsu, serta distribusi aliran peristiwa melalui Apache Kafka dan Apache Spark Structured Streaming menuju \textit{data mart} PostgreSQL yang divisualisasikan pada dasbor Grafana. Pada evaluasi eksperimental, model deteksi mencapai \textit{Mean Average Precision} (mAP@50) sebesar 91,4\%; pemrosesan tepi beroperasi stabil pada \textit{throughput} 18,9 FPS dan latensi inferensi 53 milidetik; sedangkan pialang pesan mencatat tingkat keberhasilan pengiriman 100\% dengan median latensi 2,65 milidetik. Pengujian operasional pada data lapangan selama 10 jam (07.00--17.00 WIB) menyaring 11.172 deteksi mentah menjadi 113 kejadian pelanggaran aktual. Dasbor memetakan distribusi kepadatan bimodal dengan puncak anomali tercatat pada pukul 14.45--15.15 WIB yang berkorelasi dengan faktor cuaca dan dinamika pergantian kelas. Infrastruktur terintegrasi tersebut menggeser pengawasan reaktif menjadi instrumen Sistem Pendukung Keputusan (\textit{Decision Support System}) yang proaktif dalam perencanaan patroli keselamatan di lingkungan kampus.

\noindent
\textbf{Kata kunci:} ByteTrack, \textit{Data Mart}, Pelanggaran Helm, YOLOv8 % masukkan keyword
\end{abstractind}
